\documentclass[11pt]{article}
\usepackage[margin=1in]{geometry}
\usepackage{hyperref}
\pagestyle{empty}
\setlength{\parindent}{0pt}\setlength{\parskip}{0.7em}
\begin{document}

26 October 2026

The Editors\\
\emph{Empirical Software Engineering}\\
Springer Nature

Dear Editors,

I submit for your consideration the manuscript \textbf{``What Should a Graph Node
Say? A Controlled Study of Node Descriptions for Cross-Project Vulnerability
Detection''} as a research article.

\textbf{What the paper does.} Graph-based vulnerability detectors report good
accuracy on projects seen in training and much worse accuracy on projects they
have not seen. I ask how much of that failure depends on what each node of the
program graph tells the network. With the graph, network and training protocol
held fixed, I compare four node descriptions (identifier text, text with
identifiers renamed, the parser's grammar kind, and a hand-built operation
label) on uniform random samples of BigVul and DiverseVul, using
project-disjoint folds, a cluster bootstrap over held-out projects, direct
paired tests, and controls for function size and for single dominant projects.

\textbf{Main findings, including the negative ones.} Text embeddings learned
from scratch transfer close to chance. Renaming identifiers or using grammar
kinds improves AUC substantially, and the two do about equally well. The
hand-built operation label adds nothing to the grammar kind, which I explain
with an exact information-theoretic argument. Much of the gain disappears when
function size is held fixed, and a one-line function-size rule has a higher mean
AUC than every network, while on BigVul a classifier that sees only formatting
(trailing whitespace is present in 61\% of vulnerable and 2\% of benign
functions) reaches AUC 0.896. A fine-tuned CodeBERT falls from 0.922 on raw
source to 0.728 once formatting is removed; it then beats the graph models by a
small margin but is not distinguishable from the size rule. I report these results at face value and
recommend that studies of cross-project detectors include a function-size
baseline and a vocabulary-free baseline.

\textbf{Fit with the journal.} The paper is an empirical study with a strong
emphasis on evaluation methodology: controlled comparison, explicit threats to
validity, annotation agreement for the labelling function, and every number
generated by released scripts from stored predictions. It falls under
``verification and validation'' and ``applications of artificial intelligence
techniques to software engineering'' in the journal's aims and scope.

\textbf{Transparency.} An earlier version of this work, with a different
(federated-learning) framing, was submitted to another journal and not
accepted; a second version with the present representation framing was also
reviewed elsewhere and not accepted. The present manuscript responds to those
reviews with new experiments (random samples, ten matched folds, additional
baselines, a Transformer and a pretrained model, ablations, statistical tests,
function-size controls and a re-done annotation protocol) and withdraws the
claims that the new evidence did not support. The manuscript is not under
consideration elsewhere. All code, annotation sheets and per-fold predictions
are available at \url{https://github.com/vinhqdang/vulmorph-fed}. Limitations
that I could not remove (a short fine-tuning schedule for the pretrained model, two language-model
annotators and no human second annotator, samples of about 10{,}000 functions)
are stated in the paper.

Thank you for your consideration.

Yours sincerely,

Quang-Vinh Dang\\
British University Vietnam, Hung Yen, Vietnam\\
vinh.dq4@buv.edu.vn\\
ORCID: 0000-0002-3877-8024

\end{document}
